\documentclass[12pt]{article}
% Préambule commun à tous les documents extraits de « La Revue CAPES »
\usepackage[T1]{fontenc}
\usepackage[utf8]{inputenc}
\usepackage{lmodern}
\usepackage{amsmath,amssymb,amsthm,amsfonts,mathrsfs}
\usepackage{setspace}
\usepackage{listings}
\usepackage{pgfplots}
\pgfplotsset{compat=1.18}
\usepackage{longtable}
\usepackage{adjustbox}
\usepackage{lettrine}
\usepackage{tkz-tab}
\usepackage{changepage}
\usepackage{pdflscape}
\usepackage{graphicx}
\usepackage{tikz}
\usepackage{xcolor}
\usepackage[a4paper,top=2cm,bottom=2.5cm,left=2.5cm,right=2cm]{geometry}
\usepackage{fancyhdr}
\usepackage{draftwatermark}
\SetWatermarkText{La RevueCapes}
\SetWatermarkLightness{0.9}
\SetWatermarkScale{2}
\usepackage{hyperref}
\hypersetup{colorlinks=true,linkcolor=blue!50!black,urlcolor=blue!50!black}

\definecolor{revue}{RGB}{20,60,120}

\pagestyle{fancy}
\fancyhf{}
\renewcommand{\headrulewidth}{0.4pt}
\setlength{\headheight}{15pt}

% \entete{Matière}{Titre}{Type de document}
\newcommand{\entete}[3]{%
  \fancyhead[L]{\small\textsc{La Revue CAPES} -- #1}%
  \fancyhead[R]{\small #2}%
  \fancyfoot[C]{\thepage}%
  \thispagestyle{plain}%
  \begin{center}
    {\color{revue}\rule{\textwidth}{1.2pt}}\\[4pt]
    {\small\textsc{Concours directs CAP-CEG / CAPES -- Burkina Faso}}\\[6pt]
    {\LARGE\bfseries\color{revue} #2}\\[6pt]
    {\large\itshape #1 -- #3}\\[2pt]
    {\color{revue}\rule{\textwidth}{1.2pt}}
  \end{center}
  \vspace{0.5cm}%
}

% Encadré pour les remarques ajoutées dans les corrigés
\newcommand{\remarque}[1]{\par\medskip\noindent\fcolorbox{revue}{revue!6}{\parbox{\dimexpr\linewidth-2\fboxsep-2\fboxrule}{\textbf{\color{revue}Remarque.} #1}}\par\medskip}


\begin{document}
\entete{Mathématiques}{Sujet type n°5}{Énoncé}
\begin{spacing}{1.5}

 
 \textbf{Exercice 1}
 
 1) Soit $A$ la matrice carrée : 
$A = \begin{pmatrix}
5 & -6 \\
1 & 0
\end{pmatrix}$

 1.a) Trouver les valeurs propres $\lambda_1$ et $\lambda_2 (\lambda_1 < \lambda_2)$ et des vecteurs propres $v_1$ et $v_2$ associés de la matrice $A$.  

On notera $D$ la matrice diagonale formée par les valeurs propres de $A$ rangées dans 
l'ordre croissant. 

b) Donner l'expression d'une matrice régulière $P$ telle que $A = PDP^{-1}$.

c) En déduire la matrice $A^n$, $n$ étant un entier naturel strictement positif. 

2. On considère la suite récurrente d'ordre 2 définie par : 

$U_{n+2} = 5U_{n+1}-6U_n$
avec $U_0 = U_1 = 1$

On note $V_{n+2}$ le vecteur colonne $\begin{pmatrix}
U_{n+2} \\ 
U_{n+1}
\end{pmatrix} $

a) Montrer que $V_{n+2} = AV_{n+1}$

b) En déduire l'expression du terme général $U_n$ en fonction de $n$. 

c) Calculer \[\lim_{n\longrightarrow +\infty}U_n\]
 
\medskip\textbf{Exercice 2}
 
 Soit $\theta$ un réel appartenant à l'intervalle $J = ]-\frac{\pi}{2}, + \frac{\pi}{2}[$. 

On considère l'équation $(E)$ de la variable complexe $z$ : 

$(E): z^2\cos^2\theta-4z\cos\theta + 5 - \cos^2\theta = 0$ 

1. Résoudre $(E)$ dans l'ensemble des complexes. On précisera pour quelle(s) valeur(s) de $\theta$ l'équation $(E)$ admet une racine double et la valeur de cette racine. 

2. Le plan complexe étant rapporté à un repère orthogonal, on note par $M_1$ et $M_2$ les 
points du plan complexe dont les affixes respectives sont $z_1$ et $z_2$, solutions de $(E)$. 

Donner l'équation cartésienne de la courbe du plan, lieu géométrique de $M_1$ et $M_2$ lorsque $\theta$ varie dans l'intervalle $J$.

\medskip\textbf{Exercice 3 }
 
 1. Calculer la dérivée de la fonction $f :x\longmapsto (x^2+1)\sin x$,
 
 2. Montrer de deux manières différentes que l'équation $(x^2+1)\cos x+2x\sin x=0$ admet au moins une solution dans $[0,\pi]$,
 
3. Soient $a,b,c\in \mathbb{R}$. En utilisant le théorème de Rolle,démontrer qu'il existe $x\in ]0,1[$, tel que
 $4ax^3+3bx^2+2cx=a+b+c$
 
4. En utilisant le théorème des accroissements finis à la fonction $f:x\longmapsto \ln x$, démontrer que $\forall {x>0}, \frac{x}{x+1}<\ln(x+1)<x$.

\medskip\textbf{Exercice 4}
 
  1. Les fonctions suivantes définies pour $x\in \mathbb{R}$ par:
 
 $i)f(x)=\frac{1-\cos x}{x^2}$ 
 
 $ii)g(x)= \frac{1}{x}$ 
 
 sont-elles prolongeables par continuité en 0?
 
2. Soit la fonction définie par : $f(x)=\begin{cases}
\ln (x+1), \hspace{0,5 cm} si \hspace{0,5 cm} x\geq 0\\
\sin{x}, \hspace{0,5 cm} si\hspace{0,5 cm} x<0
\end{cases}$

 (a) Étudier la continuité, la dérivabilité et la continuité éventuelle de la dérivée de la fonction $f$ sur son domaine de définition $D_f$.
 
 (b) $f$ est-elle de classe $C^1$ sur $D_f$?
 
 \quad
 

\end{spacing}
\end{document}
